\begin{thanks}

本论文的研究工作是在导师的悉心指导下完成的。从选题的确定、研究问题的梳理，到技术路线的论证、数值实验的设计和论文的反复修改，导师都投入了大量时间。研究推进过程中，导师不断提醒我先说明问题边界，再给出能够由数据或推导支持的结论；“结论要有可核对的依据”也成为我组织公式、参数、计算与讨论的基本要求。导师对论证层次和学术表达的严格要求，使我在修改中逐步学会区分研究假设、计算结果与结论的适用范围。在此向导师致以诚挚的谢意。

感谢课题组各位老师和同学。在每周的讨论中，大家指出了本文模型中若干含混的假设，并帮助我把物理量、边界条件和量纲约定整理清楚。围绕状态变量如何定义、不同误差来源怎样区分、评价指标能够支持何种判断，课题组的交流使我有机会从不同角度检查自己的推理。特别感谢在数值实验阶段提供帮助的同学，他们协助完成了网格收敛性验证与参数扫描的部分算例，并与我一起核对计算步骤和误差指标。讨论中形成的认真核对、及时反馈的习惯，也帮助我更有条理地推进后续工作。

感谢实验室的工程师和实验人员。他们在设备调试、数据采集和误差标定方面提供了专业支持，使本文在设计算例时能够考虑测量流程、仪器误差与数据处理之间的联系，而不只停留在纸面推导。感谢实验室提供的计算条件，使参数扫描和重采样计算能够按计划开展。每一次对采集条件和处理过程的说明，都提醒我计算结果必须结合数据来源与测量边界来理解；这些具体而细致的帮助，对本文的研究与写作都有重要意义。

感谢参与论文评审和预答辩的各位专家。各位专家关于不确定度报告方式、评价指标选择和结论表述边界的意见，使我重新检查了论证中证据与判断之间的对应关系。根据这些意见，论文删去了若干支撑不足的表述，并补充了区间估计、误差分解和适用范围的说明。专家们严谨、具体的建议不仅帮助我完善了论文，也让我认识到清楚交代研究限制与报告计算结果同样重要。

感谢家人一直以来的理解和支持。在论文工作较为紧张的阶段，他们对我的体谅和照顾使我能够安心完成研究与写作；在进展不顺或需要反复修改时，他们给予的鼓励也让我保持耐心。家人的支持使我能够在长期、细致的工作中继续前行，这份理解我会一直铭记。

最后，感谢所有在本课题研究过程中给予帮助的老师、同学和工作人员，也感谢国家自然科学基金等项目及相关单位提供的支持。谨以此文向所有关心和帮助过我的人表示诚挚谢意。论文仍有不足，文中的疏漏与责任由作者本人承担。

\end{thanks}
